\documentclass[aspectratio=169]{beamer}
\usetheme{default}
\usepackage{booktabs}
\definecolor{editblue}{RGB}{38,117,168}
\setbeamercolor{structure}{fg=editblue}
\title{Can action reconstruction make an editor obey?}
\subtitle{Token-aligned sequence-edit JEPA: LDAD coefficient screen}
\date{17 July 2026}
\begin{document}
\frame{\titlepage}

\begin{frame}{Low prediction error is not enough}
  \begin{itemize}
    \item The model sees a damaged solution and a literal insert, delete, or replace action.
    \item It should predict the state produced by that particular edit.
    \item An action-blind model can predict an average nearby state and still obtain low error.
    \item We therefore swap the action at evaluation; a causal predictor should become at least 5\% worse.
  \end{itemize}
  \vfill
  \textbf{Question:} does reconstructing the action from the observed latent change force the forward model to use its action?
\end{frame}

\begin{frame}{Latent Difference Action Decoding adds an auxiliary task}
  \begin{itemize}
    \item LDAD means Latent Difference Action Decoding.
    \item A decoder receives only the online latent difference between consecutive observed states.
    \item It reconstructs the complete observed edit phrase, including operation, position, and content.
    \item Five matched cells test LDAD weights 1, 10, and 20 and two representation-regularization controls.
    \item All runs use seed 0, 2,000 curriculum trajectories, three epochs, zero dropout, and 256 held-out examples.
  \end{itemize}
\end{frame}

\begin{frame}{LDAD improves dynamics but misses the causal gate}
  \centering
  \begin{tabular}{lrrrr}
    \toprule
    Condition & One-step & Recursive & Shuf./match & Rank \\
    \midrule
    No VICReg, LDAD 10 & .078 & .140 & 1.007 & 137.6 \\
    VICReg .02, LDAD 1 & .106 & .190 & 1.010 & 131.2 \\
    VICReg .02, LDAD 10 & .078 & .140 & 1.007 & 137.5 \\
    VICReg .02, LDAD 20 & \textbf{.076} & \textbf{.134} & 1.006 & \textbf{137.9} \\
    VICReg .1, LDAD 10 & .089 & .159 & \textbf{1.011} & 136.8 \\
    \bottomrule
  \end{tabular}
  \vfill
  \begin{itemize}
    \item Every model beats persistence and reconstructs about 82\% of action tokens.
    \item No shuffled-action ratio approaches the required 1.05 threshold.
    \item LDAD 20 has the best errors but a peak pre-clipping gradient norm of 241.7.
  \end{itemize}
\end{frame}

\begin{frame}{What to notice and why it matters}
  \begin{itemize}
    \item The auxiliary decoder can read actions from observed state changes while the forward predictor still underuses the supplied action.
    \item LDAD is provisionally useful as a dynamics regularizer, not as a causal-action solution.
    \item Mask-only action sensitivity disappears when that checkpoint is evaluated on mixed edits, exposing a distribution confound.
    \item Non-curriculum runs also repeated corruptions despite a configuration requesting fresh examples; this is now fixed and tested.
    \item Next: matched fixed/fresh data controls, common cross-regime evaluation, and one lower-learning-rate LDAD-20 control.
  \end{itemize}
  \vfill
  \textbf{Gate:} no hierarchy or planning claim until mixed-edit action sensitivity reaches 1.05.
\end{frame}
\end{document}
